\subsubsection{Matrix Translation of Projected-State Pairings}
\label{sec:matrix-state-pairings}
The intrinsic fixed-projector formulas in
Eqs.\ \eqref{eq:fixed-projector-amplitudes}--\eqref{eq:projected-block-density}
have the following ordinary-matrix translation. The operation
\(\Hop\) is represented by the ordinary Hermitian operation, as established
in Eq.\ \eqref{eq:matrix-operation-dictionary}.

The fixed projector and lowering element are represented by
\begin{equation*}
\pv{\Pi}(+)
\ \longleftrightarrow\
\smat{1&0\\0&0},
\qquad
\pv{a}(-)
\ \longleftrightarrow\
\smat{0&0\\1&0}.
\end{equation*}
Consequently, the two intrinsic expansions used in
Eq.\ \eqref{eq:hermitian-column-pairing} are represented by
\begin{equation*}
\begin{aligned}
\pv{\psi}_{\eqtext{block}}
=a\pv{\Pi}(+)+b\pv{a}(-)
&\ \longleftrightarrow\
\matf{\pv{\psi}_{\eqtext{block}}}=\smat{a&0\\b&0},\\
\pv{\phi}_{\eqtext{block}}
=c\pv{\Pi}(+)+d\pv{a}(-)
&\ \longleftrightarrow\
\matf{\pv{\phi}_{\eqtext{block}}}=\smat{c&0\\d&0}.
\end{aligned}
\end{equation*}
The intrinsic pairing is therefore translated into the familiar
row--column product by
\begin{equation*}
\begin{aligned}
2\,\eqfunc{scalar}\!\left(
\herm{\pv{\phi}_{\eqtext{block}}}\pv{\psi}_{\eqtext{block}}
\right)
&=\eqfunc{tr}\!\left(
\mherm{\matf{\pv{\phi}_{\eqtext{block}}}}
\matf{\pv{\psi}_{\eqtext{block}}}\right)\\
&=\mherm{\mcol{c,d}}\,\mcol{a,b}
=c^*a+d^*b.
\end{aligned}
\end{equation*}
This is the ordinary-matrix realization of
Eq.\ \eqref{eq:hermitian-column-pairing}. Algebraic scalar extraction is
thereby identified as one half of the ordinary trace in this realization.

For the nonzero state scale \(N\) from
Eq.\ \eqref{eq:projected-block-scale}, the opposite-order product is
represented by
\begin{equation*}
\matf{\frac{
\pv{\psi}_{\eqtext{block}}
\herm{\pv{\psi}_{\eqtext{block}}}}{N}}
=\frac{1}{N}
\mcol{a,b}\mrow{a^*,b^*}
=\frac{1}{N}\smat{|a|^2&a b^*\\b a^*&|b|^2}.
\end{equation*}
The displayed matrix is Hermitian, is unchanged when squared, has determinant
zero, and has exactly one independent column direction. These are the
ordinary-matrix statements
corresponding to the intrinsic identities in
Eq.\ \eqref{eq:projected-block-density}. The name ``pure state'' refers to
this single normalized state rather than to a statistical mixture
\cite{sakurai2017}.

An unrestricted complex paravector is represented by an arbitrary
two-by-two complex matrix. Its two columns are acted upon independently by
left multiplication. Under right multiplication by
\(\pv{\Pi}(+)\), the first column is retained and the second
is set to zero, producing the matrices displayed above. If that
right-projector restriction is omitted, two independent copies of each
left-acting Pauli or Dirac equation are carried by the two columns.

The ordered pair
\(\bPsi_{\eqtext{Dirac}}
=(\bPsi_{\eqtext{Dirac},1},\bPsi_{\eqtext{Dirac},2})\)
defined intrinsically above
can also be displayed as the block column
\begin{equation*}
\bPsi_{\eqtext{Dirac}}
\ \longleftrightarrow\
\begin{bmatrix}
\bPsi_{\eqtext{Dirac},1}\\[2pt]
\bPsi_{\eqtext{Dirac},2}
\end{bmatrix}.
\end{equation*}
Likewise, the component action of the general pair operator may be packaged
as the block matrix
\begin{equation*}
\mat{\mathcal O}
:=
\begin{bmatrix}
\pv{\mathcal O}(1,1)&\pv{\mathcal O}(1,2)\\
\pv{\mathcal O}(2,1)&\pv{\mathcal O}(2,2)
\end{bmatrix}.
\end{equation*}
With this optional packaging, the intrinsic operator pairing from
Eq.\ \eqref{eq:dirac-operator-pairing} becomes
\begin{equation*}
\begin{aligned}
&\left\langle\bPsi_{\eqtext{Dirac\_left}}\middle|
\mat{\mathcal O}\bPsi_{\eqtext{Dirac\_right}}\right\rangle\\
&=2\iiint\,\eqfunc{scalar}\!\Big(
\herm{\bPsi_{\eqtext{Dirac\_left},1}}
\pv{\mathcal O}(1,1)\bPsi_{\eqtext{Dirac\_right},1}\\
&\qquad
+\herm{\bPsi_{\eqtext{Dirac\_left},1}}
\pv{\mathcal O}(1,2)\bPsi_{\eqtext{Dirac\_right},2}\\
&\qquad
+\herm{\bPsi_{\eqtext{Dirac\_left},2}}
\pv{\mathcal O}(2,1)\bPsi_{\eqtext{Dirac\_right},1}\\
&\qquad
+\herm{\bPsi_{\eqtext{Dirac\_left},2}}
\pv{\mathcal O}(2,2)\bPsi_{\eqtext{Dirac\_right},2}
\Big)\,\dd r_1\,\dd r_2\,\dd r_3.
\end{aligned}
\end{equation*}
No additional state components or physical assumptions are introduced by
either matrix display; both are translations of the ordered-pair and
fixed-projector algebra defined above.
